% Part VII: the sentiment-modulated risk-budgeting pipeline (nabia_python/nabia/sentiment_rb).

\chapter{Overview of the pipeline}\label{ch:rb-overview}

The pipeline implements the specification \code{doc/sentiment\_risk\_budgeting\_spec.md}: each portfolio asset
(equity, commodity, government bond, credit) is mapped to a set of sentiment indicators, inherited down an
indicator tree; an elastic net per asset aggregates them into one sentiment score; the score is normalized to
$[0,1]$ and modulates the asset's risk budget, or, in long-short mode, decides whether the asset is held long or
short; a risk-budgeting optimizer turns budgets into weights; the result is backtested against plain equal risk
contribution (ERC). News sentiment comes from a language model. Every numerical step that the C++ library
implements runs in the \code{nabia\_python\_bridge} wheel:
\begin{center}\small
\begin{tabularx}{\linewidth}{@{}l X l@{}}
\toprule
Step & Method & nabia \\
\midrule
News aggregation, volatility, smoothing & exponentially weighted averages & \code{signal.exponential\_moving\_average} \\
Elastic net per asset & coordinate descent, $(\alpha,\lambda)$ grid, purged walk-forward CV, block bootstrap & \code{RegressionEngine} (\code{frames.regress}) \\
Normalization (option) & $\Phi(S/\hat\sigma_S)$ & \code{stats.normal\_cdf} \\
Weights & risk budgeting (Spinu), risk contributions & \code{risk.RiskBudgeting}, \code{risk.Volatility} \\
Long-only backtest, ERC benchmark & walk-forward engine with a budget schedule & \code{frames.backtest} \\
Metrics & Sharpe, drawdown, turnover, \dots & \code{strategy.compute\_performance} \\
\bottomrule
\end{tabularx}
\end{center}
pandas and numpy do the bookkeeping (alignment, rolling z-scores, ranks, correlations). The long-short backtest
runs on a small simulator that repeats the engine's accounting with a direction per asset
(\S\ref{sec:rb-simulator}).

\paragraph{Data flow.} For one configuration, \code{run} (\coderef{rb_run}) executes:
\begin{enumerate}
  \item synthetic point-in-time market data, indicators and news (Chapter~\ref{ch:rb-synthetic});
  \item the news pipeline: dedup, triage, scoring, daily digest, node features (Chapter~\ref{ch:rb-news});
  \item point-in-time transforms and the candidate features of each asset (Chapter~\ref{ch:rb-features});
  \item the walk-forward elastic net, the score $S$ and its normalization $s$ (Chapter~\ref{ch:rb-model});
  \item budgets or directions, the weights and the backtests (Chapter~\ref{ch:rb-portfolio});
  \item tables, CSV outputs and figures.
\end{enumerate}
\snippet{rb_run}{Python}{The pipeline for one configuration}

\paragraph{Code layout} (\code{nabia\_python/nabia/sentiment\_rb}):
\begin{center}\small
\begin{tabularx}{\linewidth}{@{}l l X@{}}
\toprule
Folder & Files & Area \\
\midrule
--- & \code{config.py}, \code{pipeline.py}, \code{\_\_main\_\_.py} & settings, \code{run()}, command line and sweep \\
\code{data/} & \code{synthetic.py}, \code{real.py}, \code{example\_news.py} & synthetic data; real prices, indicators and news files; example news \\
\code{features/} & \code{tree.py}, \code{transforms.py} & indicator tree; point-in-time transforms, candidate features, target \\
\code{llm/} & \code{schemas.py}, \code{client.py}, \code{news.py}, \code{lexicon.py}, \code{batch.py} & schemas; mock or cached Anthropic client; dedup, scoring, node features; offline reader; Message Batches \\
\code{model/} & \code{elastic\_net.py} & elastic net per asset, walk-forward $S$ and $s$ \\
\code{portfolio/} & \code{budgets.py}, \code{backtest.py} & budgets and directions; backtests, simulator, checks \\
\code{report/} & \code{tables.py}, \code{figures.py} & report tables and CSV; figures \\
\code{prompts/} & \code{triage\_v1.md}, \dots & versioned prompt files (system and user templates) \\
\bottomrule
\end{tabularx}
\end{center}
All settings are in one dictionary, \code{CONFIG} (\coderef{rb_config}); \code{load\_config}
(\coderef{rb_load_config}) deep-merges a JSON file of overrides, so taxonomy, features, horizons, bounds, models
and the normalization method change without code edits.

%==============================================================================
\chapter{Indicator tree, point-in-time features and target}\label{ch:rb-features}

\section{The indicator tree}
The taxonomy is a tree of nodes: \code{L0.GLOBAL} (global risk appetite) at the root; the asset classes
\code{L1.EQ}, \code{L1.CMD}, \code{L1.FI}; buckets \code{L2.*} (US and European equities, energy, precious
metals, 10-year Treasuries, high yield); and one node per underlying \code{L3.*}. Each node has a parent; the
ancestors of $n$ are $\anc(n)=\{n,\text{parent}(n),\dots,\code{L0.GLOBAL}\}$ (\coderef{rb_ancestors}). Asset $i$
sits at an L3 node $n_i$ and belongs to an asset class $c(i)\in\{\mathrm{EQ},\mathrm{CMD},\mathrm{GOVT},\mathrm{CREDIT}\}$.

Every feature $k$ is attached to a node $\nu(k)$ and carries a \emph{sign per asset class} $\sigma_{k,c}\in\{-1,0,1\}$
that orients it: after orientation, higher means more bullish for the asset's price. Risk-off is bullish for
government bonds, so a risk-appetite indicator has $\sigma_{k,\mathrm{GOVT}}=-1$. The candidate features of asset
$i$ are those of all its ancestors that are not switched off for its class:
\[
  \mathcal F_i=\{\,k:\ \nu(k)\in\anc(n_i),\ \sigma_{k,c(i)}\ne0\,\}.
\]
Each feature enters once (the matrix is keyed by feature), which the tests check
(\code{test\_assets\_inherit\_every\_ancestor\_feature\_once}).
\snippet{rb_asset_matrix}{Python}{Candidate features of an asset, oriented}

\section{Point-in-time alignment and transforms}
A value observed for date $t$ (\emph{as of} $t$) becomes usable $\ell_k$ days later (\emph{available at}
$t+\ell_k$): a COT report three days, financial-conditions indices five, LLM features one day after the news is
received. The usable series is therefore $\tilde x_{k,t}=x_{k,t-\ell_k}$, and nothing earlier than its
availability is ever used. The series is then made stationary and standardized:
\[
  u_{k,t}=\begin{cases}\tilde x_{k,t}-\tilde x_{k,t-5}&\text{(transform ``diff'')}\\ \tilde x_{k,t}&\text{(transform ``z'')}\end{cases},
  \qquad
  z_{k,t}=\operatorname{clip}_{[-3,3]}\!\left(\frac{u_{k,t}-m_t^{(W)}(u_k)}{s_t^{(W)}(u_k)}\right),
\]
with the rolling mean $m^{(W)}$ and standard deviation $s^{(W)}$ over $W=756$ trading days (about three years,
at least 252 observations). Rolling, not full-sample, statistics keep every value point in time. The oriented
feature of asset $i$ is $\hat z^{(i)}_{k,t}=\sigma_{k,c(i)}\,z_{k,t}$.
\snippet{rb_transform}{Python}{Point in time, stationary, rolling z-score}

\section{The target}
The elastic net predicts the forward vol-scaled return over $h=21$ days,
\[
  y_{i,t}=\frac{P_{i,t+h}/P_{i,t}-1}{\hat\sigma_{i,t}\sqrt h},\qquad
  \hat\sigma_{i,t}^2=\operatorname{EWMA}_{H=20}\big(r_{i,\cdot}^2\big)_t ,
\]
where the exponentially weighted mean of squared daily returns (half-life 20 days, smoothing factor
$\alpha=1-2^{-1/20}$) comes from \code{signal.exponential\_moving\_average}. Scaling by volatility makes targets
comparable across assets and calm or turbulent periods. $y_{i,t}$ is only known at $t+h$, which the training
windows respect (\S\ref{sec:rb-purge}).
\snippet{rb_target}{Python}{Vol-scaled forward return}
\begin{codepointers}
  \codeitem{rb_build_features}{all features: indicators, node LLM features, digest outputs}
  \codeitem{rb_node_sign}{orientation of the latent factors (synthetic data)}
\end{codepointers}

%==============================================================================
\chapter{The LLM news pipeline}\label{ch:rb-news}

\section{Prompts, schemas and the client}
Three prompts, stored as versioned files in \code{prompts/} (system and user templates, separated by a marker):
\begin{description}
  \item[A, triage] (a small model): is the article relevant, which taxonomy nodes does it concern, with which
        relevance and channel (\emph{direct}: the node's asset is the subject; \emph{spillover}: affected through
        a sector, competitor or macro link). Market recaps are marked irrelevant.
  \item[B, scoring] (a stronger model): for each node, the sentiment in $[-1,1]$ for the asset's price over the
        next 1--20 trading days, a confidence, magnitude, horizon, surprise versus expectations, fact type and a
        short rationale. Government bonds: hawkish policy and upside inflation are negative.
  \item[C, digest] (one call per day, on the day's deduplicated headlines): risk appetite, growth tone, inflation
        tone, geopolitical stress, credit tone, central-bank tone per bank, commodity supply tone per complex.
\end{description}
Each response is validated against a pydantic schema (\coderef{rb_triage_schema}, \coderef{rb_score_schema},
\coderef{rb_digest_schema}): numbers in range, categories from closed lists. A malformed response is retried up to
twice, then logged and skipped. Every call is appended to \code{llm\_audit.jsonl} with the prompt version (file
name and hash of its text), the model, a hash of the request, the article or date it concerns, the attempt, and
the raw response, so each score can be reproduced or re-scored. Templates are filled by \code{fill}
(\coderef{rb_fill}).
\snippet{rb_llm_call}{Python}{One LLM call: fill, call, validate, retry, audit}
With \code{--live}, the call goes to the Anthropic API at temperature 0, with the output forced through a tool whose
input schema is the pydantic schema, and the system prompt cached; the SDK retries rate limits and server errors
with exponential backoff. Offline, \code{mock\_response} (\coderef{rb_mock}) answers from the synthetic truth of
the article, with noise, and returns 1\% of first attempts malformed so the retry path is exercised.
\snippet{rb_llm_anthropic}{Python}{The Anthropic call: structured output through a tool}
\paragraph{Look-ahead.} A model trained after an article's date may ``know'' what followed. The pipeline offers
anonymization (entity names replaced by \emph{Entity k} in every prompt, the mapping to node IDs kept outside the
model), and limits live scoring to recent news (\code{llm.live\_days}); out-of-sample performance after the model's
training cutoff is the real validation.

\section{Deduplication and novelty}
A story is often published several times. Each headline is embedded as a hashed bag of words: tokens are
hashed into $d=256$ buckets, counted and normalized, $v_a\in\mathbb R^{256}$, $\|v_a\|=1$ (\coderef{rb_embed}).
Article $a$ joins the cluster of an earlier article $b$ received within 48 hours if $v_a^\top v_b>0.9$. Its
\emph{novelty} is
\[
  \nu_a=\frac1{1+n_a},
\]
with $n_a$ the number of earlier articles of its cluster in the window: 1 for a new story, $\tfrac12$ for the
first repeat, and so on.
\snippet{rb_dedup}{Python}{Clusters and novelty}

\section{Weights and decay-weighted aggregation}
Each (article, node) score $s_a$ gets the weight
\[
  w_a=\text{relevance}_a\times\text{confidence}_a\times\nu_a\times\omega_{\text{source}(a)},
\]
with source weights $\omega$ (wire 1, newspaper 0.8, blog 0.4). The node's sentiment on day $t$ is the
decay-weighted mean of the scores received up to $t$:
\begin{equation}\label{eq:rb-decay}
  \bar s_{n,t}=\frac{\sum_{a:\,t_a\le t}w_a\,s_a\,e^{-\lambda(t-t_a)}}{\sum_{a:\,t_a\le t}w_a\,e^{-\lambda(t-t_a)}},
  \qquad \lambda=\frac{\ln 2}{H},\quad H=3\text{ days}.
\end{equation}
\snippet{rb_weights}{Python}{Weight of each (article, node) score}
\begin{proposition}\label{prop:rb-ema}
Let $X_t=\sum_{a:\,t_a=t}w_as_a$ and $W_t=\sum_{a:\,t_a=t}w_a$ be the day sums, and $E(x)$ the exponential moving
average with smoothing factor $\alpha=1-e^{-\lambda}$ started from zero,
$E_t=\alpha x_t+(1-\alpha)E_{t-1}$, $E_{-1}=0$. Then $\bar s_{n,t}=E_t(X)/E_t(W)$ whenever $E_t(W)>0$.
\end{proposition}
\begin{proof}
Unrolling the recursion, $E_t(x)=\alpha\sum_{j\le t}(1-\alpha)^{t-j}x_j=\alpha\sum_{j\le t}e^{-\lambda(t-j)}x_j$.
Substituting the day sums, $E_t(X)=\alpha\sum_{a}w_as_ae^{-\lambda(t-t_a)}$ and
$E_t(W)=\alpha\sum_{a}w_ae^{-\lambda(t-t_a)}$; the factor $\alpha$ cancels in the ratio, which is
\eqref{eq:rb-decay}.
\end{proof}
The code therefore computes the aggregation with nabia's EWMA, prepending a zero so that it starts from $E_{-1}=0$:
\snippet{rb_ema0}{Python}{Zero-started EWMA (nabia)}

\section{The five features per node}
For every node $n$ and day $t$:
\begin{center}\small
\begin{tabular}{@{}l l l@{}}
\toprule
Feature & Definition & Orientation \\
\midrule
\code{llm\_level} & $\bar s_{n,t}$, eq.~\eqref{eq:rb-decay} & $+$ \\
\code{llm\_volume} & $\log(1+N_{n,t})$, $N$ the day's relevant articles & $-$ (attention spikes on bad news) \\
\code{llm\_dispersion} & standard deviation of the day's scores & $-$ (disagreement) \\
\code{llm\_mom5} & $\bar s_{n,t}-\bar s_{n,t-5}$ & $+$ \\
\code{llm\_surprise} & $\bar s_{n,t}-\frac1{60}\sum_{j<60}\bar s_{n,t-j}$ & $+$ \\
\bottomrule
\end{tabular}
\end{center}
LLM sentiment is defined with respect to the price of the node's asset, so its sign is $+1$ for the node's subtree,
except at \code{L0.GLOBAL}, where it is risk appetite and takes the risk-on sign of each class. Like every indicator
they are lagged (one day), z-scored and oriented (Chapter~\ref{ch:rb-features}).
\snippet{rb_node_features}{Python}{The five LLM features of every node}

\section{The daily digest}
Prompt C's outputs become features of their node, with a sign per class, e.g. inflation tone negative for equities,
bonds and credit but positive for commodities, a hawkish Fed negative for everything, energy supply tightening
positive for the energy bucket (\code{digest\_features} in \coderef{rb_config}). Null outputs (no relevant news for
a central bank) count as neutral.
\begin{codepointers}
  \codeitem{rb_run_news}{the news pipeline: dedup, triage, scoring, digests}
\end{codepointers}

%==============================================================================
\chapter{Elastic net per asset and the sentiment score}\label{ch:rb-model}

\section{The model}
For each asset $i$ and refit date $T$, the elastic net (Chapter~\ref{ch:en}) is fitted on the oriented
candidates $\hat z^{(i)}_{k,t}$, $k\in\mathcal F_i$:
\[
  \min_{\beta_0,\beta}\ \frac1{2n}\sum_{t\in\mathcal T_T}\big(y_{i,t}-\beta_0-\textstyle\sum_k\beta_k\hat z^{(i)}_{k,t}\big)^2
  +\lambda\Big(\alpha\|\beta\|_1+\tfrac{1-\alpha}2\|\beta\|_2^2\Big),
\]
with standardized features, $\alpha\in\{0.1,0.3,0.5,0.7,0.9\}$ and a $\lambda$ path, chosen by cross-validation.
The fit is nabia's regression engine (\code{frames.regress}) with all candidates in one subset, so the engine tunes
only $(\alpha,\lambda)$.

\section{Purged walk-forward training}\label{sec:rb-purge}
Targets overlap: $y_{i,t}$ spans $t+1,\dots,t+h$. Two precautions keep the future out:
\begin{enumerate}
  \item \textbf{The training rows end $h$ days before the refit date}: $\mathcal T_T=\{t\le T-h\}$, the rows whose
        target is known at $T$ (an assertion checks it). The window is expanding (or rolling, by configuration),
        and starts once enough history exists for the z-scores; a missing z-score counts as neutral (0).
  \item \textbf{Walk-forward cross-validation with a gap}: the engine's folds train on the past and validate on
        the next block, with $h$ rows removed between them (\code{validation\_gap}), the analogue of
        \code{TimeSeriesSplit(gap=h)}.
\end{enumerate}
Refits happen at the end of each quarter (or month); the assets of one refit date are fitted in parallel threads,
since the engine releases Python's GIL.
\snippet{rb_fit_at}{Python}{Training rows of one refit: purged and point in time}

\section{Positivity and stability selection}
The features are pre-oriented, so a negative coefficient would be a nonsensical sign flip; the specification asks
for $\beta\ge0$. nabia's elastic net has no sign constraint, so the pipeline applies it by elimination: fit, drop
the features with $\beta_k<0$, refit, until all coefficients are non-negative. At the same time it applies
\emph{stability selection}: each fit includes a moving-block bootstrap of $B=50$ samples (blocks of $h$ rows, so
overlapping targets stay together), and the selection frequency
\[
  \pi_k=\frac1B\sum_{b=1}^B\one\big[\beta^{(b)}_k\ne0\big]
\]
must exceed $0.6$; features below are dropped and the model refitted. If nothing survives, $\beta=0$: the asset's
score is neutral.
\snippet{rb_fit_asset}{Python}{Elastic net with positivity and stability selection}

\section{The score and its normalization}
With the coefficients of the latest refit, the sentiment score is the prediction minus the intercept,
\[
  S_{i,t}=\sum_{k\in\mathcal F_i}\beta_{ik}\,\hat z^{(i)}_{k,t},
\]
so all-zero coefficients give $S=0$. It is normalized to $[0,1]$ without full-sample statistics, which would look
ahead and be dominated by outliers:
\begin{description}
  \item[Rolling percentile] (default): the share of the last $L=756$ scores (computed with the current
        coefficients) below today's, counting ties as half,
        \[
          s^{\mathrm{ts}}_{i,t}=\frac1L\sum_{u=t-L}^{t-1}\Big(\one[S_{i,u}<S_{i,t}]+\tfrac12\,\one[S_{i,u}=S_{i,t}]\Big).
        \]
        The mid-rank makes a constant score (an asset without features) exactly neutral, $s=\tfrac12$.
  \item[Normal CDF]: $s^{\mathrm{ts}}_{i,t}=\Phi\big(S_{i,t}/\hat\sigma_{S,t}\big)$, with the rolling standard
        deviation $\hat\sigma_S$ (\code{stats.normal\_cdf}).
\end{description}
An optional cross-sectional rank across assets,
$s^{\mathrm{xs}}_{i,t}=(\operatorname{rank}_i S_{\cdot,t}-1)/(N-1)$, is blended,
$s=\alpha_{\mathrm{ts}}s^{\mathrm{ts}}+(1-\alpha_{\mathrm{ts}})s^{\mathrm{xs}}$ ($\alpha_{\mathrm{ts}}=1$ by default),
and the result is smoothed by an EWMA with a 5-day half-life. Every value is checked to lie in $[0,1]$.
\snippet{rb_normalize}{Python}{Score and rolling percentile with the current coefficients}
\snippet{rb_blend}{Python}{Cross-sectional blend and EWMA smoothing}

\section{In-sample and out-of-sample diagnostics}
For every refit, the pipeline records the information coefficient (IC, the correlation of the score with the
realized target) on the training rows and on the next, unseen block. In-sample IC above out-of-sample IC measures
overfitting; the report also gives the IC of $S$ and of $s$ per asset and pooled over the whole out-of-sample
period, and the pooled IC of each LLM feature.
\begin{codepointers}
  \codeitem{rb_walk_forward}{walk-forward refits, scores, normalization, IC records}
  \codeitem{rb_tables}{report tables: comparison, IC, snapshot, selection frequencies}
\end{codepointers}

%==============================================================================
\chapter{From sentiment to positions}\label{ch:rb-portfolio}

\section{Risk budgeting}
Given a covariance $\Sigma$ and budgets $b_i>0$, $\sum_ib_i=1$, the risk-budgeting portfolio
(Chapter~\ref{c:sec:riskbudgeting}) solves
\[
  \min_{w>0}\ R(w)-\sum_ib_i\ln w_i,\qquad R(w)=\sqrt{w^\top\Sigma w},
\]
and is rescaled to sum to one. At the optimum the relative risk contributions equal the budgets:
\[
  \mathrm{RC}_i(w)=\frac{w_i(\Sigma w)_i}{w^\top\Sigma w}=b_i .
\]
(The specification's $\tfrac12w^\top\Sigma w$ gives the same weights up to scale.) The pipeline checks this at
every rebalance with nabia's \code{RiskBudgeting} to $10^{-4}$; observed errors are below $10^{-8}$
(\coderef{rb_check_rc}).

\section{Long-only: modulated budgets}
With base budgets $b^0$ (ERC weights per asset) and normalized sentiment $s_i$,
\[
  b_i=\frac{b_i^0\,\big(m_{\min}+(m_{\max}-m_{\min})\,s_i\big)}{\sum_jb_j^0\,\big(m_{\min}+(m_{\max}-m_{\min})\,s_j\big)},
  \qquad m_{\min}=0.5,\ m_{\max}=1.5,
\]
so neutral sentiment ($s=\tfrac12$) keeps $b^0$ and the multiplier ranges over $[0.5,1.5]$. Budgets stay strictly
positive (the log barrier needs $b>0$). The change per rebalance is capped at $\pm25\%$ relative: the budgets are
clipped to $[b^{\mathrm{prev}}(1-c),\,b^{\mathrm{prev}}(1+c)]$ and renormalized, repeatedly.
\snippet{rb_modulate}{Python}{Budget modulation with a change cap}

\section{Long-short: directions}\label{sec:rb-longshort}
In long-short mode, sentiment decides the side of each position and the budgets stay $b^0$ (optionally scaled by
conviction, $s_i$ for a long and $1-s_i$ for a short). The direction $d_i\in\{+1,-1\}$ follows a neutral band with
hysteresis, so an asset hovering around one threshold does not flip back and forth:
\[
  d_{i,t}=\begin{cases}-1&s_{i,t}<s_{\mathrm{short}}\ (0.35)\\ +1&s_{i,t}>s_{\mathrm{long}}\ (0.5)\\ d_{i,t^-}&\text{otherwise.}\end{cases}
\]
\snippet{rb_directions}{Python}{Directions with hysteresis}
A short position earns $-r_i$. With $D=\operatorname{diag}(d)$, the directed returns are $Dr$, whose covariance is
$\Sigma'=D\Sigma D$: a shorted asset's correlations with the others change sign, so it is sized as the hedge or
diversifier it becomes. The optimizer runs on $\Sigma'$ and returns positive exposures $w'$; the weights are
$w=Dw'$.
\begin{proposition}\label{prop:rb-short}
Let $w'>0$ be the risk-budgeting solution on $\Sigma'=D\Sigma D$ with budgets $b$, and $w=Dw'$. Then the risk
contributions of $w$ on the original covariance equal the budgets: $\mathrm{RC}_i(w;\Sigma)=b_i$. Every position,
long or short, carries its risk budget.
\end{proposition}
\begin{proof}
Since $d_i^2=1$, $w^\top\Sigma w=w'^\top D\Sigma Dw'=w'^\top\Sigma'w'$ and
$w_i(\Sigma w)_i=d_iw'_i(\Sigma Dw')_i=w'_i\,(D\Sigma Dw')_i=w'_i(\Sigma'w')_i$. Hence
$\mathrm{RC}_i(w;\Sigma)=\mathrm{RC}_i(w';\Sigma')=b_i$ by optimality of $w'$.
\end{proof}
\snippet{rb_solve}{Python}{Risk budgeting on $D\Sigma D$ (nabia)}

\section{The long-short simulator}\label{sec:rb-simulator}
nabia's backtest engine has no notion of direction, and directions change between rebalances while the covariance
window spans several of them, so flipping columns of its input would mix signs within one window. The simulator
\code{simulate} repeats the engine's accounting (Chapter~\ref{c:sec:engine}) and applies the current directions to
the whole trailing window at each rebalance:
\begin{itemize}
  \item \textbf{Covariance}, the engine's estimator: weights $\omega_j\propto2^{-a_j/H}$ for a row of age $a_j$
        (half-life $H=60$; equal weights for the sample estimator), normalized, and
        \[
          \hat\Sigma=\frac{\sum_j\omega_j(r_j-\bar r)(r_j-\bar r)^\top}{1-\sum_j\omega_j^2},\qquad\bar r=\sum_j\omega_jr_j ,
        \]
        over $W=252$ rows (\coderef{rb_window_cov}).
  \item \textbf{Rebalances} at rows $W-1+kf$, $f=21$; the budgets and directions in force are the last schedule row
        at or before the date.
  \item \textbf{Drift}: between rebalances each holding grows with its return,
        $w\leftarrow w\circ(1+r_t)/(1+w^\top r_t)$; short weights drift the same way.
  \item \textbf{Costs}: at a rebalance the turnover is $\tau=\|w^\ast-w\|_1$ and the cost $c\,\tau$ ($c=5$ bp); the
        net return is $(1+w^\top r_t)(1-c\tau)-1$.
  \item \textbf{Metrics} from the first trade on, with nabia's \code{compute\_performance}.
\end{itemize}
Run with all-long directions, the simulator reproduces nabia's engine: equity to $6.5\times10^{-10}$, weights to
$6\times10^{-9}$, identical Sharpe ratio. The long-only path asserts this on every run, and a test checks it on
independent data (\coderef{rb_test_simulator}); \coderef{rb_test_short} checks Proposition~\ref{prop:rb-short}.
\snippet{rb_simulate}{Python}{The engine's accounting with directions}
\snippet{rb_backtests}{Python}{Benchmark and overlay: same covariance, rebalancing and costs}

%==============================================================================
\chapter{Real data and live scoring}\label{ch:rb-real}

With \code{--data real} (\code{"data": \{"source": "real"\}}), \code{run} replaces the synthetic generators by three
loaders and leaves everything downstream unchanged.

\section{Prices and indicators}
\code{real\_market} downloads daily closes from Yahoo Finance and series from FRED and caches each as a CSV file in
\code{data.cache\_dir}, so a run is reproducible and works offline once the cache exists. The mapping is in
\code{CONFIG["real"]}: an asset ticker per asset (the energy asset uses Brent, since WTI settled negative on 2020-04-20)
and, per indicator, its source:
\begin{center}\small
\begin{tabular}{@{}l l l@{}}
\toprule
Feature & Source & Series \\
\midrule
\code{VIX}, \code{MOVE}, \code{USD} & Yahoo Finance & \code{\^{}VIX}, \code{\^{}MOVE}, \code{DX-Y.NYB} \\
\code{HY\_OAS}, \code{NFCI}, \code{TERM\_PREMIUM} & FRED & \code{BAMLH0A0HYM2}, \code{NFCI}, \code{THREEFYTP10} \\
\code{CRACK} & derived & $(2\,\text{gasoline}+\text{heating oil})\times42/3-\text{crude}$, in \$/barrel \\
\bottomrule
\end{tabular}
\end{center}
Features without a free source are removed from the run. Returns use the US trading calendar, other markets'
holidays forward-filled (at most three days); each indicator is taken as of its last observation on or before each
date, and the publication lags of \code{CONFIG["features"]} then apply as usual. FRED publishes only about three years
of the high-yield OAS; before that the feature is missing and counts as neutral.
\snippet{rb_real_market}{Python}{Real returns and indicator levels}

\section{News files and the offline reader}
\code{load\_news} reads a JSON file of articles with \code{publish\_ts}, \code{received\_ts} (ISO 8601), \code{source},
\code{headline} and \code{body}: the fields a news feed provides. Until one exists, \code{data/example\_news.json} stands
in: \code{example\_news.py} applies the oracle construction of \S\ref{sec:rb-oracle} to the \emph{real} assets' forward
returns (IC 0.2 by default) and writes articles in the real format, so a real-data run has news with a known skill. It
looks ahead on purpose and must be replaced by real news for any conclusion.

Offline, articles without a synthetic truth are read by \code{lexicon\_response}: the node is the taxonomy node whose
name or alias (\code{CONFIG["aliases"]}) is the longest found in the headline, its parent a spillover; the tone is
$\tanh$ of the summed polarities of finance words; recaps are recognized by their wording; the digest averages the day's
tones per node. On the example vocabulary this recovers the five tone levels in order.
\snippet{rb_lexicon_response}{Python}{The offline reading of real-format articles}

\section{Live scoring, the response cache and Message Batches}
With \code{--live}, the days from \code{llm.live\_days} before the end are scored by the model; \code{0} scores every
day. Each call has a key, the hash of the prompt version, the model and the filled request; a valid response is
appended to \code{llm.cache} (\code{data/llm\_cache.jsonl}) and every later call with the same key, in this run or the
next, is served from it, so each article is paid for once and a run can be repeated exactly. With \code{--batch}, the
uncached requests of a run are first submitted as Message Batches, one per stage, because scoring depends on triage and
the digests on the day's articles: all triage requests, then the scoring requests of the relevant articles, then the
digests. Each batch is polled until it ends, its valid responses are cached, and the pipeline's calls that follow are
cache hits; failed entries are retried one by one by the normal path.
\snippet{rb_batch_submit}{Python}{Message Batches: submit, wait, cache}
\begin{codepointers}
  \codeitem{rb_load_news}{news files}
  \codeitem{rb_make_example_news}{example news from real returns}
  \codeitem{rb_llm_request}{cache key of a call}
  \codeitem{rb_llm_prefill}{batch backfill of one stage}
\end{codepointers}

\section{Downloading news: nabia.news}\label{sec:rb-news-download}
The package \code{nabia.news} collects real articles into the JSON format above (\code{python -m nabia.news fetch}):
\begin{description}
  \item[RSS / Atom feeds] (\coderef{rb_news_feeds}): financial media and central banks, parsed with the standard library
        (RFC 822 or ISO 8601 dates, HTML stripped). Feeds carry only recent items, so \code{received\_ts} is the download
        time and history is built by running the download daily.
  \item[GDELT raw files] (\coderef{rb_news_gkg}): GDELT publishes every 15 minutes a file of the English web articles it
        processed, with their page titles. For history, two slices per business day (14:30 and 20:30 UTC) are read in
        memory; titles about the portfolio's topics are kept, established financial outlets first, at most
        \code{--per-day} per day; \code{received\_ts} is the slice's timestamp, when GDELT had the article, so the
        history stays point in time. GDELT's search API (\code{--gdelt}) is also supported but rate-limits hard.
  \item[The store] (\coderef{rb_news_store}): \code{data/news.json}, deduplicated by URL and by normalized headline on the
        same day, appended in place and written atomically; each record keeps its feed and URL.
\end{description}

\section{Scoring with a local model and scores files}
Besides the Anthropic API, the LLM steps can run on a local model served by Ollama (\code{"provider": "ollama"};
\coderef{rb_llm_ollama}): the same prompts, the output constrained to the pydantic schema through Ollama's structured
output, temperature 0, and the same response cache. On the development machine (Apple silicon, 34 GB),
\code{gemma4:26b}, a mixture-of-experts model, was the fastest model with sensible answers: about 85 generated tokens
per second and about 2 s per call, against 12--22 s for a dense 27B model. Serving several requests at once
(\code{OLLAMA\_NUM\_PARALLEL} $>1$) slowed every request several times on that machine (prompts re-read, generation
at a few tokens per second per slot), so a local model is used with one worker.

\code{python -m nabia.news score} runs triage, scoring and the daily digests over a news store and writes a
\emph{scores file} (\coderef{rb_write_scores}): per article its triage and node scores (sentiment, confidence,
relevance, novelty and the resulting weight), per day its digest, with the provider, model and prompt versions. With
\code{--scores} (\code{data.scores\_file}), the pipeline reads it (\coderef{rb_load_scores}), maps every article and
digest to the first trading date at or after its day, and calls no model.

%==============================================================================
\chapter{Synthetic data, experiments and results}\label{ch:rb-synthetic}

\section{Synthetic market data}
So that the pipeline runs offline and its behaviour can be tested against a known truth, the data are synthetic
(\coderef{rb_synthetic_market}). Every taxonomy node $n$ has a latent factor, a standardized AR(1),
$f_{n,t}=\phi f_{n,t-1}+\sqrt{1-\phi^2}\,\varepsilon_{n,t}$, $\phi=0.97$ (half-life about three weeks). Asset $i$'s
daily return is
\[
  r_{i,t}=\frac{\sigma_i}{\sqrt{252}}\Big(\kappa\,\frac{\sum_{n\in\anc(n_i)}\sigma^{\mathrm{f}}_{n,c(i)}f_{n,t}}{\sqrt{|\anc(n_i)|}}
  +\ell_{c(i)}\,g_t+\sqrt{1-\ell_{c(i)}^2}\,\eta_{i,t}\Big)+\mathrm{SR}_{c(i)}\frac{\sigma_i}{252},
\]
with annual volatilities $\sigma_i$, signal strength $\kappa=0.15$, a common shock $g$ with class loadings $\ell$
(government bonds negative), and a risk premium of Sharpe ratio $\mathrm{SR}$ per class; the factor sign
$\sigma^{\mathrm f}$ is the risk-on sign at \code{L0.GLOBAL} and $+1$ below. Each indicator is a noisy reading of
its node's factor with a slowly drifting level (or, for transform ``diff'', a stock whose changes carry the signal);
some indicators are pure noise, for the model to discard.

\section{Synthetic news with a known skill}\label{sec:rb-oracle}
The default news generator is an \emph{oracle} with a set skill (\coderef{rb_news_truth}). For node $n$, let
$x_{n,t}$ be the mean over its assets of the oriented forward target, shifted by the LLM lag so that it starts when
the news becomes usable, and $z_{n,t}$ its standardized value. The node's true tone is
\[
  T_{n,t}=\rho\,z_{n,t}+\sqrt{1-\rho^2}\,\epsilon_{n,t},\qquad\operatorname{corr}(T_{n,t},z_{n,t})=\rho ,
\]
where $\rho$ is the \emph{news IC} (\code{synthetic.sentiment\_ic}, 0.2 by default): 0 is pure noise, 1 perfect
foresight. About four articles a day are drawn (more often for global and asset-class nodes), each with tone
$\tanh(T_{n,t}+0.3\,\epsilon_a)$; 12\% are market recaps (irrelevant) and 10\% re-publications (duplicates); the mock
LLM adds scoring noise; the digest derives from the same truths. The rest of the pipeline is unchanged: it must find
the edge through point-in-time features, the purged elastic net, normalization and the budget rules. Setting
\code{"sentiment": "latent"} ties the news to the latent factors instead.
\snippet{rb_news_truth}{Python}{Oracle news truth}

\paragraph{The information coefficient.} The IC of a signal is its correlation with the subsequent realized return,
here $\operatorname{corr}(S_{i,t},y_{i,t})$, per asset or pooled over assets. Real signals have ICs of a few percent;
by the fundamental law of active management the information ratio grows roughly like
$\mathrm{IC}\times\sqrt{\text{breadth}}$, so small ICs applied often can matter.

\section{Results}
The figures and tables below come from runs with the default configuration (seed 7); \code{doc/build.sh --figures}
regenerates them. Signals start in 2018 after the warm-up of the z-scores and the first training window.

\pipelinefigure{lo_portfolio.png}{Long-only overlay against plain ERC (news IC 0.2): cumulative PnL net of costs,
rolling one-year Sharpe ratio, drawdown, return by year. Shaded: the covariance warm-up before the first
out-of-sample signal.}{fig:rb-lo-portfolio}
\pipelinefigure{lo_budgets.png}{Long-only: modulated budgets against $b^0$, normalized sentiment, weights and
cumulative risk contributions (which track the budgets).}{fig:rb-lo-budgets}
\pipelinefigure{lo_model.png}{Elastic net diagnostics: in-sample against out-of-sample IC per asset and per refit,
and the share of refits in which each feature has a non-zero coefficient.}{fig:rb-lo-model}

\IfFileExists{figures/comparison_long_only.tex}{\begin{table}[htbp]\centering\small
\caption{Long-only overlay against plain ERC (news IC 0.2).}\label{tab:rb-lo}\begin{tabular}{@{}lrr@{}}
\toprule
metric & ERC ($b^0$) & sentiment overlay \\
\midrule
annualized return & -0.0178 & -0.0123 \\
annualized volatility & 0.0503 & 0.0506 \\
sharpe ratio & -0.3313 & -0.2181 \\
max drawdown & 0.2688 & 0.2455 \\
annualized turnover & 0.6737 & 1.0579 \\
total costs & 0.0024 & 0.0038 \\
\bottomrule
\end{tabular}

\end{table}}{}
\IfFileExists{figures/comparison_long_short.tex}{\begin{table}[htbp]\centering\small
\caption{Long-short overlay against plain ERC (news IC 0.2).}\label{tab:rb-ls}\begin{tabular}{@{}lrr@{}}
\toprule
metric & ERC ($b^0$) & sentiment overlay \\
\midrule
annualized return & 0.0448 & 0.0803 \\
annualized volatility & 0.0488 & 0.0445 \\
sharpe ratio & 0.9231 & 1.7574 \\
max drawdown & 0.1236 & 0.0876 \\
annualized turnover & 0.6505 & 5.3629 \\
total costs & 0.0023 & 0.0194 \\
\bottomrule
\end{tabular}

\end{table}}{}

\paragraph{Does correct sentiment produce alpha?} The sweep reruns the pipeline on the same market data for several
news ICs. Overlay Sharpe ratio and excess return over ERC grow with the news IC; below an IC of about 0.1 the
elastic net mostly cannot separate the news from noise, and the market indicators alone give a small edge.
\IfFileExists{figures/sweep_long_only.tex}{\begin{table}[htbp]\centering\small
\caption{News-quality sweep, long-only (ERC Sharpe ratio for reference).}\label{tab:rb-sweep-lo}
\begin{tabular}{@{}lrrrrr@{}}
\toprule
news IC & ERC Sharpe & overlay Sharpe & excess return & IC of $S$ & IC of LLM level \\
\midrule
0.0 & -0.33 & -0.30 & +0.14\% & 0.12 & -0.00 \\
0.05 & -0.33 & -0.30 & +0.14\% & 0.13 & 0.04 \\
0.1 & -0.33 & -0.29 & +0.19\% & 0.16 & 0.08 \\
0.2 & -0.33 & -0.22 & +0.55\% & 0.29 & 0.15 \\
0.3 & -0.33 & -0.14 & +0.94\% & 0.43 & 0.23 \\
0.5 & -0.33 & -0.04 & +1.45\% & 0.62 & 0.37 \\
\bottomrule
\end{tabular}
\end{table}}{}
\IfFileExists{figures/sweep_long_short.tex}{\begin{table}[htbp]\centering\small
\caption{News-quality sweep, long-short.}\label{tab:rb-sweep-ls}
\begin{tabular}{@{}lrrrrr@{}}
\toprule
news IC & ERC Sharpe & overlay Sharpe & excess return & IC of $S$ & IC of LLM level \\
\midrule
0.0 & -0.33 & 0.04 & +1.87\% & 0.12 & -0.00 \\
0.05 & -0.33 & 0.13 & +2.26\% & 0.13 & 0.04 \\
0.1 & -0.33 & 0.60 & +4.73\% & 0.16 & 0.08 \\
0.2 & -0.33 & 1.57 & +9.58\% & 0.29 & 0.15 \\
0.3 & -0.33 & 2.69 & +16.23\% & 0.43 & 0.23 \\
0.5 & -0.33 & 2.91 & +18.08\% & 0.62 & 0.37 \\
\bottomrule
\end{tabular}
\end{table}}{}
\pipelinefigure{sweep_long_only.png}{Long-only sweep: Sharpe ratio, excess return and realized IC against the
news IC.}{fig:rb-sweep-lo}
\pipelinefigure{sweep_long_short.png}{Long-short sweep.}{fig:rb-sweep-ls}
\pipelinefigure{ls_budgets.png}{Long-short: budgets stay $b^0$; the sentiment panel shades the neutral band;
negative weights are short positions.}{fig:rb-ls-budgets}

\paragraph{Reading the results.} Long-only modulation can at most halve or increase by half a budget, capped at
25\% per rebalance, so even good news moves the allocation modestly. Long-short acts on the direction of the view
and captures much more of it, at the price of turnover (several times a year) and costs; with uninformative news it
earns almost nothing, which also shows that no future information leaks through the pipeline. The ERC benchmark's
negative return is a property of this seed's synthetic premiums; only the difference matters.

\section{Running and extending}
Chapter~\ref{ch:guide} lists the commands. The command line (\coderef{rb_main}) runs one configuration or a sweep;
figures are produced by \coderef{rb_plots} and \coderef{rb_plot_sweep}. To use real data, replace
\code{synthetic\_market} and \code{synthetic\_news} (\coderef{rb_synthetic_news}) by loaders returning the same
frames: returns by asset, indicator levels by \emph{as-of} date, and articles with publish and received
timestamps; everything downstream is unchanged.
